% status: ZERO
% Chapter 3 of the second edition. Plan: docs/STRUCTURE.md. Rules: AUTHORING.md.
% Measurements: research/measurements/2026-09-23-m1-llama32-3b.md
\chapter{The Roofline: Why Decode Leaves Your GPU Idle}\label{ch:roofline}

\begin{ChapterIntent}
A laptop generating text with a 3-billion-parameter model pulls nearly as many
bytes per second from memory as its GPU can stream, while its arithmetic units
do a twentieth of what they could. This chapter explains why with the roofline
model: arithmetic intensity, the two ceilings of memory bandwidth and peak
compute, and the ridge point where they meet. Decode, at a few FLOPs per byte,
sits far down the bandwidth slope; prefill, at hundreds or thousands, sits
under the compute roof. Batch size is the lever that moves decode along the
slope, and the chapter measures it --- including where the lever sticks. It
ends with what the roofline leaves out, because the best kernels in later
chapters live in that gap.
\end{ChapterIntent}

\OpeningSection{Nearly all of the memory bus, a twentieth of the arithmetic}

Take an ordinary laptop: a 2020 MacBook Pro with Apple's M1 chip, 16~GB of
unified memory and an 8-core GPU. Load Llama~3.2~3B, quantized to about four
and a half bits per weight, into llama.cpp and ask it to generate. On the day
this chapter was drafted, on a busy desktop with browsers and terminals open,
it produced \textbf{21.0 tokens per second} (measured: 21.02~$\pm$~0.30 over
five runs of 128 tokens; the record is in
\texttt{research/measurements/2026-09-23-m1-llama32-3b.md}).

Now count bytes. The model file holds 2,011,539,712 bytes of tensors, and ---
as \Chref{forward-pass} shows --- generating one token reads essentially all
of them once: this model ties its input and output embeddings, so even the
323-megabyte vocabulary table is read in full by the output projection at
every step. Twenty-one tokens per second times 2.01~GB per token is
\textbf{42~GB/s} of weights streamed out of memory.

How much could the memory deliver? The M1's memory is nominally about
68~GB/s (derived from LPDDR4X-4266 on a 128-bit bus \todo{verify the memory
type from a primary source}; consistent with Apple's statement that the M1
Max's 400~GB/s is ``nearly 6x'' the M1's). No
program gets the nominal figure; the useful ceiling is what a plain streaming
kernel achieves on the same GPU at the same moment. We measured it: 46--50~GB/s
for large reads and copies. Decode is running at \textbf{85--92\% of the
bandwidth a streaming kernel can reach}.

Now count arithmetic. A forward pass costs about two floating-point operations
per parameter per token --- one multiply and one add for every weight
(\cite{kaplan2020scaling}, eq.~2.2) --- so each token costs about 6.4~GFLOP,
and 21 tokens per second is 135~GFLOP/s. Apple quotes the M1 GPU at 2.6
TFLOP/s. Decode is using \textbf{about 5\% of the arithmetic}.

The fast part of the chip is waiting for the slow part. That is not a flaw in
llama.cpp or a quirk of Apple silicon: the napkin math below shows the same
imbalance, larger, on data-center GPUs, and a discrete consumer card shows it
too \todo{measure: re-run on an RTX 5060 Ti (448 GB/s); a prior session
reported 85.3 tok/s for a 4.7 GB Q4 7B model, about 89\% of nominal
bandwidth}. It is a property of the workload. This chapter explains which property, and what
changes it.

\NapkinSection{Napkin math: bytes per token, tokens per second}

If every token must stream $W$ bytes of weights through a memory system of
bandwidth $\beta$, then no implementation can generate faster than
\begin{equation}
r_{\max} = \frac{\beta}{W}\quad[\text{tokens/s}],
\label{eq:decode-ceiling}
\end{equation}
for a single sequence, whatever its arithmetic units do. For the M1 run:
$68.3/2.01 = 34$ tokens/s against the nominal bandwidth, $49/2.01 = 24$ against
the achieved one, and 21 measured. The model predicts the measurement to
within 15\% from two numbers, before any code is read.

The same arithmetic prices data-center hardware. Qwen3-8B has 8.19~billion
parameters (derived from its published configuration). In BF16, two bytes per
weight, it is 16.4~GB; batch-one decode reads about 15.1~GB of it per token,
because its $151{,}936\times4{,}096$ input embedding table, 1.24~GB, is
gathered one row at a time rather than streamed. An H100 SXM delivers
3.35~TB/s, so $r_{\max}\approx 3.35/15.1 \approx 220$ tokens/s. A B200,
at 8~TB/s, gives about 530. Store the same weights in FP8 and both numbers
double; in four bits they double again. Nothing about the tensor cores ---
989 dense BF16 TFLOP/s on the H100, 2.25~PFLOP/s on the B200 --- appears in
these predictions (all figures from the vendors' datasheets, logged in
\texttt{research/FRONTIER.md}). Decode speed at batch one is a memory
bandwidth question.

The arithmetic tells the other half. One token of an 8-billion-parameter model
is about 16~GFLOP. At 220 tokens/s that is 3.6~TFLOP/s --- \textbf{0.4\%} of
the H100's dense BF16 peak.

\section{Arithmetic intensity}

The ratio that decides which resource binds is the \Term{arithmetic intensity}
of a computation: the floating-point operations it performs per byte it must
move from memory,
\begin{equation}
I = \frac{F}{Q}\quad[\text{FLOP/byte}],
\label{eq:intensity}
\end{equation}
where $F$ counts operations and $Q$ counts bytes crossing the boundary we care
about --- here, between off-chip memory (DRAM, HBM or GDDR) and the chip. $Q$
is the \emph{compulsory} traffic: every byte the computation must read or write
at least once. Real kernels move more; the intensity of a kernel can only be
lower than that of the computation it implements.

For a decode step, the dominant term is a matrix--vector product: each weight
is read once and used once, in one multiply and one add. With $b_w$ bytes per
weight,
\begin{equation}
I_{\text{decode}} \approx \frac{2}{b_w}.
\label{eq:intensity-decode}
\end{equation}
BF16 weights give 1~FLOP/byte; the M1 run's weights average 0.63 bytes per
parameter, giving 3.2~FLOP/byte. Either way the number is small, and it is a
property of the product itself: a vector touches each weight once.
\Appref{matmul} derives the same result for the scalar kernels of the
book's teaching engine.

\section{Two roofs and a ridge}

A processor offers two rates: peak compute $\pi$ (FLOP/s) and memory bandwidth
$\beta$ (bytes/s). A kernel with intensity $I$ can go no faster than either
allows. This is the \Term{roofline} model of Williams, Waterman and Patterson
\cite{williams2009roofline}:
\begin{equation}
P \le \min\left(\pi,\; \beta I\right)\quad[\text{FLOP/s}].
\label{eq:roofline}
\end{equation}
Plotted on logarithmic axes, the bound is a sloped line --- the bandwidth roof,
$P=\beta I$ --- that meets a flat line --- the compute roof, $P=\pi$ --- at the
\Term{ridge point}
\begin{equation}
I^{*} = \frac{\pi}{\beta}.
\label{eq:ridge}
\end{equation}
Left of the ridge a kernel is memory-bound: faster arithmetic would not help.
Right of it, compute-bound: faster memory would not help.

\EngineFigure{ch03-roofline-m1}{The roofline of an Apple M1 GPU with five
measured llama.cpp points for Llama~3.2~3B Q4\_K\_M. Batch-one decode sits on
the bandwidth slope at 3.2~FLOP/byte; batching moves it right; prefill of 512
tokens sits under the compute roof. Intensities use compulsory bytes (weights
read once per step); the roofs are Apple's compute figure and our measured and
nominal bandwidth. Measured on a busy desktop; one machine.}{fig:ch03-roofline-m1}

The ridge is the single most useful number about a machine for inference
engineering. The table below computes it for the machines this book uses
(derived from the specifications in \texttt{research/FRONTIER.md}; dense
peaks).

\begin{center}
\begin{tabular}{@{}lrrr@{}}
\toprule
Machine & $\beta$ & $\pi$ & Ridge $I^{*}$ \\
\midrule
Apple M1 GPU & 68 GB/s & 2.6 TFLOP/s & 38 FLOP/B \\
NVIDIA H100 SXM, BF16 & 3.35 TB/s & 989 TFLOP/s & 295 FLOP/B \\
NVIDIA H100 SXM, FP8 & 3.35 TB/s & 1,979 TFLOP/s & 591 FLOP/B \\
NVIDIA H200, BF16 & 4.8 TB/s & 989 TFLOP/s & 206 FLOP/B \\
NVIDIA B200, BF16 & 8 TB/s & 2.25 PFLOP/s & 281 FLOP/B \\
NVIDIA B200, FP4 & 8 TB/s & 9 PFLOP/s & 1,125 FLOP/B \\
\bottomrule
\end{tabular}
\end{center}

\section{Decode sits on the slope}

Compare the decode intensity with the ridges. At 1--3~FLOP/byte, batch-one
decode is one to three orders of magnitude to the left of every ridge in the
table. The fraction of peak compute it can use is at most $I/I^{*}$: 8\% on the
M1 at 3.2~FLOP/byte, 0.3\% on an H100 running BF16 weights. The measured 5\% on
the M1 is right where the model says it must be.

This is the fact the rest of the book is built on. Shazeer made the same
observation about incremental decoding in 2019 and designed multi-query
attention to reduce the bytes it moves \cite{shazeer2019mqa}; Pope and
colleagues built their analysis of large-model inference on it
\cite{pope2023scaling}. Every technique that follows either raises the
intensity of decode, reduces the bytes it must move, or accepts the bound and
hides it.

\section{Batch size moves the point}

The simplest way to raise intensity is to reuse each weight for more than one
token. If $B$ sequences are decoded together, one pass over the weights serves
all of them: the FLOPs grow by $B$ and the weight bytes do not, so
\begin{equation}
I_{\text{decode}}(B) \approx \frac{2B}{b_w},
\qquad
B^{*} \approx \frac{b_w\, I^{*}}{2}.
\label{eq:batch-ridge}
\end{equation}
Below $B^{*}$, adding a sequence costs almost nothing: the step time is set by
the weight bytes, and throughput grows linearly with $B$. Above $B^{*}$, the
step becomes compute-bound and its time grows with $B$. For BF16 weights on an
H100, $B^{*}\approx 295$; for four-bit weights computed in BF16, about 74; on
the M1 with its 0.63-byte weights, about 12--17.

We measured it. Decoding 1 to 32 sequences of Llama~3.2~3B together on the
M1, the aggregate rate rose from 20 to 176 tokens per second: \textbf{8.8$\times$
the tokens for 3.7$\times$ the step time} (measured, same record). An ideal
kernel would have delivered 32$\times$ the tokens for about twice the step
time (derived: at $B=32$ a step needs 206~GFLOP, 79~ms at the M1's peak,
against 41~ms to stream the weights once at 49~GB/s). The measurement follows
the model in aggregate and departs from it in detail, and the detail is
instructive.

\EngineFigure{ch03-batch-staircase}{Measured step time of batched decode on
the M1 as the number of sequences grows. The roofline predicts a flat line
until the ridge; the measurement is a staircase set by which kernel ggml
selects for each batch size. Shaded regions are the kernel regimes read from
ggml's source at llama.cpp commit \texttt{389ff61d}; the measured binary is
build \texttt{d00685831}.}{fig:ch03-batch-staircase}

Two and three sequences cost almost as much as two and three separate steps.
Four and five are cheaper than three. From twelve to thirty-two the step time
is nearly flat, but at three times the batch-one time. None of this is in
equation~\eqref{eq:roofline}. It is in the kernels, and it is where this
chapter's lab and its case study go.

\section{Prefill sits under the flat roof}

A prompt of $S$ tokens is processed in one forward pass. Every token needs the
same weights, so one pass over the weights serves all $S$: prefill is a batch
of size $S$ in equation~\eqref{eq:batch-ridge}. For a 512-token prompt on the
M1, the intensity is about $512\times 3.2\approx 1{,}600$~FLOP/byte, forty
times the ridge. Prefill should be compute-bound, and it is: the M1 processed
the prompt at \textbf{248 tokens per second}, about 1.6~TFLOP/s or 61\% of the
GPU's peak --- twelve times the per-token speed of decode on the same weights
(measured). \Chref{prefill-decode} takes this difference apart.

\section{Rooflines of real machines}

Three trends show up as soon as the table above is drawn to scale.

First, ridges are high and rising. Data-center GPUs need hundreds of FLOPs per
byte before their tensor cores are the constraint, and each new precision
doubles the compute roof without moving the bandwidth roof: the B200's FP4
ridge is about 1,100~FLOP/byte. Low precision raises decode's intensity (fewer
bytes per weight) but raises the ridge in step, so the gap stays.

Second, bandwidth grows more slowly than compute. From B200 to B300 NVIDIA's
listed dense FP4 throughput rises by half while per-GPU HBM bandwidth stays at
8~TB/s (NVIDIA HGX and GB300 NVL72 specifications, logged in
\texttt{research/FRONTIER.md}). The ridge moves right again.

Third, the practical roofs sit below the datasheet. On the M1 the streaming
kernel reached 70\% of nominal bandwidth; the engine reached 85--92\% of that.
\todo{measure: the same three numbers --- nominal, streaming-kernel and decode
bandwidth --- on an RTX 5060 Ti and an H100}.

\FigureIntent{Rooflines of the M1, RTX 5060 Ti, H100 and B200 on one set of
log--log axes, with each machine's batch-one decode point for an 8B model in
BF16 and in four bits: the gap between the decode points and the ridges is
the chapter's argument in one picture.}

\section{What the roofline leaves out}

The roofline is a bound, not a prediction. The staircase in
Figure~\ref{fig:ch03-batch-staircase} shows four things it does not model,
each of which later chapters take up.

\emph{Tile quantization.} A matrix kernel works in fixed tiles; if its tile is
32 batch columns wide, nine sequences pay for thirty-two. \Chref{matmul-hardware}.

\emph{Occupancy.} A kernel must give the processor enough independent work to
hide memory latency. Decode attention, with one query per sequence, often
cannot, and then neither roof binds. \Chref{decode-attention}.

\emph{Fixed per-step costs.} Launching kernels, synchronizing and scheduling
take time that neither FLOPs nor bytes explain; at batch one they can be a
large share of a step. \Chref{kernel-languages}.

\emph{Bytes that do not batch.} Batching amortizes weights, not the KV cache:
every sequence reads its own cache every step, so at long context the
attention part of decode stays on the bandwidth slope however large the batch.
\Chref{kv-cache}.

\LedgerSection

Batching is the first technique in the book, and the roofline lets us charge
it precisely.

\begin{ConstraintLedger}
Compute & buys & Utilization rises about $B$-fold until the ridge: on the M1, 5\% of peak at $B=1$ to about 43\% at $B=32$ (measured). \\
Memory bandwidth & moves & Weight bytes are shared across the batch; KV bytes are not (\Chref{kv-cache}). \\
Latency & spends & Each sequence waits for the whole batch: 3.7$\times$ the step time at $B=32$ on the M1 (measured). \\
Memory capacity & spends & Every sequence in the batch holds its own KV cache. \\
Correctness & risks & Results can change with batch size unless kernels are batch-invariant (\Chref{fast-wrong-answers}). \\
\end{ConstraintLedger}

\LabSection{draw your own roofline}

\LabIntent{In \texttt{code/labs/ch03-roofline/}, a dependency-free Rust
program that (1) measures streaming read and copy bandwidth with 1--$n$
threads, (2) measures a blocked matrix product's FLOP/s for $B = 1, 2, 4,
\dots, 512$ columns against a fixed weight matrix, and (3) writes both as CSV
for a plotting script. The reader predicts $B^{*}$ from equation
\eqref{eq:batch-ridge} before running it, then explains every point that falls
below both roofs. An oracle compares every product with a naive
implementation before anything is timed. Extension: run the same sweep with
the scalar kernels of \Appref{matmul} and explain why they sit far below both
roofs.}

\HappenedSection{a kernel below both roofs}

The staircase is not noise; it is a design decision recorded in ggml's source.
At llama.cpp commit \texttt{389ff61d} (2026-05-10), the Metal backend chooses
among three matrix kernels for K-quant weights by the number of batch columns
\todo{verify: confirm the same selection in the measured build
\texttt{d00685831}}:
\begin{itemize}
\item one to three columns: a matrix--vector kernel that handles one column per
pass over the weights --- so two sequences read the weights twice;
\item four to eight: a small-batch kernel that handles three, four or five
columns per pass, chosen by a table in the source;
\item more than eight: a tiled matrix--matrix kernel whose tiles, on GPUs
without Metal's newer tensor instructions, span 64 weight rows and 32 batch
columns.
\end{itemize}
Each kernel is right for its regime and the thresholds are reasonable; the
source comments that the parameters are not known to be optimal across
hardware. The lesson is general. A roofline says batching is nearly free until
the ridge. Whether it is free on your machine depends on whether the kernel
you are actually running reuses each byte it loads --- and that is a property
of code, not of the workload.

\FrontierSection{2026-09-23}

\emph{Asymmetric scaling.} The FlashAttention-4 authors describe Blackwell as a
machine on which ``tensor core throughput doubles while other functional units
scale more slowly'' and redesign attention around the units that did not
scale \cite{zadouri2026flashattention4}. The roofline's single compute roof is
becoming several: tensor-core math, other arithmetic, special functions,
shared-memory bandwidth.

\emph{The roof keeps rising faster than the slope.} B300-class parts add FP4
throughput without adding HBM bandwidth per GPU. Unless model architectures
cut bytes per token --- the subject of Parts~\ref{part:one-request}
and~\ref{part:architectures} --- the gap between decode and the ridge grows.

\todo{verify: announced next-generation (Rubin-class) HBM bandwidth and FP4
throughput, from NVIDIA's own specifications, before stating a trend.}

\ExercisesSection

\begin{enumerate}
\item \emph{Prediction.} A 70-billion-parameter model stored in FP8 runs on
one H200. Predict its batch-one decode rate from
equation~\eqref{eq:decode-ceiling}. Which parameters did you have to assume?
\item \emph{Derivation.} Derive $B^{*}$ in equation~\eqref{eq:batch-ridge}
for weights stored in four bits but multiplied in BF16. Why does the answer
depend on where dequantization happens?
\item \emph{Implementation.} Run the chapter's lab and report your machine's
streaming bandwidth, peak matrix FLOP/s and ridge. How far is your measured
$B^{*}$ from the predicted one?
\item \emph{Falsification.} Find a configuration of llama.cpp (or any engine)
where batch-one decode is more than twice as slow as
equation~\eqref{eq:decode-ceiling} predicts. Identify the cause with a
profiler before guessing.
\end{enumerate}
\todo{worked solutions for exercises 1 and 2 in \texttt{tex/worked/ch03-roofline.tex}}
